\begin{resumo}[Abstract]
\begin{otherlanguage*}{english}

This study aimed to develop a web-based system to support psychopedagogical assessment in early childhood education, bringing together learner registration, management of assessment instruments, recording of instrument applications, and report writing in a single environment. It is an applied study conducted under the Design Science Research approach, in which the requirements were defined through a literature review, a documentary analysis of psychopedagogical instruments, and a semi-structured interview with a psychopedagogist, and the system was built incrementally. The application follows a three-tier architecture, with a React interface, a REST API in Node.js and Express, and a MySQL database accessed through Prisma. The main features include role-based and ownership-based access control, the separation between system administration and access to clinical data in compliance with professional confidentiality, automatic versioning of instruments, audit logging, and PDF report generation. All 147 automated tests passed, with 83.06\% line coverage in the backend, and the automated accessibility check found no violations of the WCAG 2.2 level A and AA criteria after the corrections. Five psychopedagogy professionals evaluated the system through a questionnaire and gave an overall mean score of 4.43 on a scale from 1 to 5, with report generation receiving the lowest mean (3.8). Of the fifteen functional requirements, thirteen were fully met and two were partially met. It is concluded that the system meets its proposed objective and was perceived as suitable for professional practice. However, it remains a prototype, since the lack of encryption at rest, of a data retention policy, and of interoperability, together with the small evaluation sample, prevents its use with real data without further adjustments.

\vspace{\onelineskip}

\noindent
\textbf{Keywords}: psychopedagogy; psychopedagogical assessment; early childhood education; web system; educational technologies.

\end{otherlanguage*}
\end{resumo}